\section{Dataset actualizado y protocolo experimental}

\subsection{Qué cambió desde la Etapa 1}

La Etapa 1 trabajó con 31\,622 imágenes. Después se añadieron ejemplos a las
clases que habían quedado más cortas: mancha gris y las tres deficiencias
nutricionales. La publicación completa consultada para esta etapa contenía
33\,437 archivos distribuidos en nueve clases \citep{huggingface2026maize}.
Todos se abrieron, se normalizó su orientación EXIF y se calculó SHA-256,
resolución y varianza del Laplaciano como indicador descriptivo de nitidez.
No hubo imágenes ilegibles. Cuatro archivos de hojas sanas repetían exactamente
el contenido de otro archivo y se eliminaron antes de dividir. El corpus
experimental quedó en 33\,433 imágenes únicas.

Esta secuencia explica cifras que antes podían parecer contradictorias.
31\,622 corresponde al corpus de la Etapa 1; 33\,437 es el número de archivos
de la publicación ampliada; 33\,433 es el total utilizado aquí después de la
deduplicación. La cifra 33\,438 que aparecía en una versión de la auditoría no
coincidió con el dataset descargado ni con su manifiesto y por eso no se
conservó como resultado.

\begin{table}[H]
  \centering
  \caption{Composición por clase antes y después de la ampliación. La fuente es
  el manifiesto generado, no una transcripción manual.}
  \label{tab:dataset-e1-e2}
  \small
  \begin{tabular}{lrrr}
\toprule
\textbf{Clase} & \textbf{Etapa 1} & \textbf{Etapa 2} & \textbf{Cambio} \\
\midrule
Roya común & 2~256 & 2~256 & +0 \\
Gusano cogollero & 4~857 & 4~857 & +0 \\
Mancha gris & 1~119 & 1~930 & +811 \\
Sana & 8~744 & 8~740 & -4 \\
Necrosis letal & 6~415 & 6~415 & +0 \\
Def. nitrógeno & 523 & 846 & +323 \\
Tizón norteño & 6~830 & 6~830 & +0 \\
Def. fósforo & 612 & 938 & +326 \\
Def. potasio & 266 & 621 & +355 \\
Total & 31~622 & 33~433 & +1811 \\
\bottomrule
\end{tabular}

\end{table}

El incremento no fue uniforme. Potasio pasó de 266 a 621 imágenes y fósforo de
612 a 938; nitrógeno llegó a 846. Mancha gris recibió la mayor ampliación
absoluta entre las clases problemáticas, de 1\,119 a 1\,930. Aun así, la clase
sana conserva 8\,740 imágenes, más de catorce veces el soporte de potasio. La
ampliación reduce el extremo del desbalance, pero no lo elimina.

\subsection{Composición y procedencia}

El manifiesto conserva catorce identificadores de fuente obtenidos del nombre
versionado de cada archivo. Las colecciones más grandes son
\texttt{maize\_africa}, \texttt{multi\_desease},
\texttt{maize\_desease\_v1.1}, \texttt{corn\_leaf\_roboflow} y
\texttt{cropdg}. No existe un identificador de finca, municipio, dispositivo o
campaña de captura. Por eso «fuente» permite detectar diferencias entre
colecciones, pero no medir generalización geográfica en El Salvador.

La etiqueta de entorno separa 3\,551 imágenes de laboratorio y 29\,882 de
campo. Su distribución está muy lejos de ser factorial: roya común concentra
2\,150 imágenes de laboratorio y 106 de campo; mancha gris tiene 513 y 1\,417;
tizón norteño 888 y 5\,942. Las otras seis clases no poseen ejemplos de
laboratorio. Esta asimetría importa más que el total global. Si el grupo de
laboratorio obtiene otro F1, la diferencia puede deberse tanto al fondo como a
las clases que lo componen.

\begin{figure}[H]
  \centering
  \includegraphics[width=0.90\textwidth]{dataset_distribution.pdf}
  \caption{Distribución del corpus único por clase y entorno. Las deficiencias
  ampliaron su soporte, pero siguen siendo minoritarias y exclusivamente de
  campo en los metadatos disponibles.}
  \label{fig:dataset-distribution}
\end{figure}

La Figura~\ref{fig:dataset-distribution} deja ver dos riesgos. El primero es el
peso de la clase sana y de necrosis letal frente a las deficiencias. El segundo
es que «laboratorio» funciona casi como una pista de clase. Esto impide
interpretar una diferencia global laboratorio/campo como efecto puro del
entorno y obliga a acompañarla con comparaciones dentro de las tres clases que
sí aparecen en ambos grupos.

\subsection{Prevención de fugas y división}

La deduplicación se hizo antes del muestreo. Cada fila contiene ruta relativa,
clase, índice de clase, entorno, fuente, ancho, alto, megapíxeles, indicador de
blur y SHA-256. El fingerprint del conjunto completo se calculó ordenando las
rutas y concatenando ruta, clase, entorno y hash. El valor resultante es:

\begin{center}
\small\ttfamily
c44de0783d3c3f88ef6151eaa75536d1020f5b208327419c4cb7f3ac459d3892
\end{center}

Se reservaron 5\,015 imágenes (15\%) como holdout con semilla 4202. El 85\%
restante se dividió en 23\,403 imágenes de entrenamiento y 5\,015 de
validación, utilizando semilla 42. Ambos cortes se estratificaron por la
combinación clase--entorno. El archivo \texttt{holdout.lock.json} se creó antes
de extraer características y guarda el fingerprint, el hash del CSV de holdout
y la regla de no consultar sus etiquetas durante selección.

El conjunto histórico llamado test no se reutilizó como prueba final. Esa
decisión sacrifica comparabilidad directa con la Etapa 1, pero recupera una
interpretación limpia: la evaluación final responde a decisiones fijadas sin
haber visto esas 5\,015 etiquetas. Los resultados de ambas etapas se comparan
más adelante como referencias sobre corpus y protocolos distintos, no como
una prueba controlada de superioridad.

\subsection{Preprocesamiento y métricas}

Las imágenes se corrigieron según EXIF, se convirtieron a RGB, se redimensionaron
a la entrada nativa de cada backbone y se normalizaron con media y desviación
de ImageNet. El redimensionamiento directo no conserva la relación de aspecto;
es una limitación deliberadamente visible del protocolo. Cada foto se decodificó
una vez durante la extracción conjunta, y FastViT recibió 256 píxeles mientras
los otros backbones recibieron 224.

La métrica de selección fue macro-F1, porque concede el mismo peso a potasio y
a la clase sana. También se guardaron accuracy, precision macro, recall macro
y F1 ponderado. Accuracy permite continuidad con la etapa anterior; F1
ponderado describe el volumen total; macro-F1 evita que cualquiera de esas dos
lecturas oculte las clases pequeñas. Para fairness se calcularon las mismas
métricas dentro de grupos, pero sus brechas se reportan como diferencias
descriptivas y no como causalidad.
